\documentclass[letterpaper,10pt]{article}
% Shared formatting for Anton Serdyuchenko's resumes.
% Compact single-column layout designed for readability and ATS-friendly text extraction.

\usepackage{fontspec}
\setmainfont{DejaVu Serif}
\usepackage{geometry}
\geometry{left=0.58in,right=0.58in,top=0.48in,bottom=0.48in}
\usepackage{enumitem}
\usepackage[hidelinks]{hyperref}
\usepackage{titlesec}
\usepackage{tabularx}
\usepackage{microtype}
\usepackage{xcolor}
\usepackage{ragged2e}

\pagestyle{empty}
\setlength{\parindent}{0pt}
\setlength{\tabcolsep}{0pt}
\urlstyle{same}
\raggedbottom
\setlength{\emergencystretch}{2em}

\titleformat{\section}
  {\vspace{-2pt}\large\bfseries\raggedright}
  {}{0em}{}
  [\titlerule\vspace{-4pt}]
\titlespacing*{\section}{0pt}{7pt}{4pt}

\newcommand{\resumeHeader}[3]{%
  \begin{center}
    {\LARGE\bfseries #1}\\[-1pt]
    {\normalsize\bfseries #2}\\[2pt]
    {\small #3}
  \end{center}
  \vspace{-7pt}
}

\newcommand{\resumeSubHeadingListStart}{%
  \begin{itemize}[leftmargin=0pt,label={},itemsep=5pt,topsep=0pt,parsep=0pt,partopsep=0pt]
}
\newcommand{\resumeSubHeadingListEnd}{\end{itemize}}

\newcommand{\resumeSubheading}[4]{%
  \item
  \begin{tabularx}{\textwidth}{X >{\raggedleft\arraybackslash}p{0.22\textwidth}}
    \textbf{#1} & {\small #2} \\
    \textit{#3} & \textit{\small #4}
  \end{tabularx}
  \vspace{-2pt}
}

\newcommand{\resumeProjectHeading}[3]{%
  \item
  \begin{tabularx}{\textwidth}{X r}
    \textbf{#1} \hspace{0.4em}{\small #2} & {\small #3}
  \end{tabularx}
  \vspace{-2pt}
}

\newcommand{\resumeItemListStart}{%
  \begin{itemize}[leftmargin=15pt,label=\textbullet,itemsep=1.5pt,topsep=2pt,parsep=0pt,partopsep=0pt]
}
\newcommand{\resumeItemListEnd}{\end{itemize}\vspace{-2pt}}
\newcommand{\resumeItem}[1]{\item {\small #1}}

\newcommand{\resumeLine}[2]{%
  \textbf{#1}: {\small #2}\\[-1pt]
}

\newcommand{\sep}{\hspace{0.45em}\textbar\hspace{0.45em}}


\begin{document}

\resumeHeader
  {Anton Serdyuchenko}
  {Software Engineer | DevOps / AI Infrastructure}
  {\href{mailto:anton415@gmail.com}{anton415@gmail.com} \sep
   \href{https://www.linkedin.com/in/antonserdyuchenko}{LinkedIn} \sep
   \href{https://github.com/anton415}{GitHub} \sep
   \href{https://serdyuchenko.com}{serdyuchenko.com}}

\section{Summary}
\small
Software engineer with 9 years of experience developing and maintaining enterprise software at the Bank of Russia. Core work includes Java and Vaadin feature development, business-logic fixes, legacy refactoring, JUnit testing, code review, SQL, and application diagnostics. Building practical DevOps and AgentOps evidence through \textit{finance-lab}, an open-source engineering lab for coding-agent workflows, CI, evaluations, workflow automation, and human approval gates.

\section{Experience}
\resumeSubHeadingListStart
  \resumeSubheading
    {Bank of Russia}{Oct 2017 -- Present}
    {Java Software Engineer (official title: Leading Engineer)}{Russia}
    \resumeItemListStart
      \resumeItem{Develop and maintain features in a long-lived enterprise application, working primarily with Java and Vaadin across business logic and UI behavior.}
      \resumeItem{Implemented a cross-layer refactoring of a core business workflow across domain logic, UI, services, and workflow orchestration, migrating affected functionality to an updated domain model and validating changes with JUnit and manual testing.}
      \resumeItem{Developed reusable Vaadin components for repeated UI patterns used across multiple application forms.}
      \resumeItem{Fix business-logic defects and refactor legacy code; use SQL and application diagnostics to investigate and verify behavior.}
      \resumeItem{Review peer changes and write or maintain JUnit tests as part of regular development and verification work.}
    \resumeItemListEnd
\resumeSubHeadingListEnd

\section{Selected Project}
\resumeSubHeadingListStart
  \resumeProjectHeading
    {finance-lab}
    {\href{https://github.com/anton415/finance-lab}{github.com/anton415/finance-lab}}
    {2026 -- Present}
    \resumeItemListStart
      \resumeItem{Created an open-source personal-finance application and engineering lab focused on reliable AI-agent-assisted software development (AgentOps).}
      \resumeItem{Established a human-gated delivery workflow with acceptance criteria, GitHub Actions CI, advisory coding-agent PR reviews, and explicit review boundaries.}
      \resumeItem{Built the project around a documented evaluation framework with 12 candidate coding-agent tasks, one manually calibrated grader, and separate application-test, agent-task, and grader-calibration evidence.}
      \resumeItem{Uses Python automation for repository workflow control and safe dry-run/live separation so issue and PR state changes remain auditable and human-approved.}
    \resumeItemListEnd
\resumeSubHeadingListEnd

\section{Skills}
\resumeLine{Software Engineering}{Java, Vaadin, SQL, JUnit, Git/GitLab, code review, debugging, legacy refactoring}
\resumeLine{Backend / Workflow}{PostgreSQL, Camunda/BPMN, Mockito, Maven}
\resumeLine{Delivery / Platform}{GitHub Actions, CI/CD, Linux, Python automation (basic)}
\resumeLine{AI Engineering}{coding-agent workflows, evaluation design, human approval gates, AgentOps experiments}

\section{Education \& Selected Training}
\resumeSubHeadingListStart
  \resumeSubheading
    {Moscow Aviation Institute (MAI)}{2022}
    {Master's degree, Management -- Economic Security, Faculty of Information Technologies}{Moscow, Russia}
\resumeSubHeadingListEnd
\small
\textbf{Selected training:} Camunda Academy -- BPMN Overview; Anthropic -- AI Fluency Framework \& Foundations; Anthropic -- Claude Code in Action.

\end{document}
